\subsection{FUNDAMENTAL DOMAIN, STABILISERS}

Recall (Integration, Chap. VII, §2, no. 10, Def. 2) that a subset D of E is called a fundamental domain for the group W if every orbit of W in E meets D in exactly one point. This is equivalent to the following pair of conditions:

\begin{enumerate}
\def\labelenumi{\alph{enumi})}
\item
  For every \(x \in \mathrm{E}\) , there exists \(w \in \mathrm{W}\) such that \(w(x) \in \mathrm{D}\) .
\item
  If \(x, y \in D\) and \(w \in W\) are such that \(y = w(x)\) , then y = x (even though we may have \(w \neq 1\) ).
\end{enumerate}

We are going to prove the following three statements:

THEOREM 2. For any chamber C, the closure \(\overline{\mathbf{C}}\) of C is a fundamental domain for the action of W on E.

PROPOSITION 1. Let \(\mathbf{F}\) be a facet and \(\mathbf{C}\) a chamber such that \(\mathbf{F} \subset \overline{\mathbf{C}}\) . Let \(w \in W\) . The following conditions are equivalent:

\begin{enumerate}
\def\labelenumi{(\roman{enumi})}
\item
  \(w(\mathbf{F})\) meets F;
\item
  \(w(\mathbf{F}) = \mathbf{F};\)
\item
  \(w(\overline{\mathrm{F}}) = \overline{\mathrm{F}};\)
\item
  w fixes at least one point of F;
\item
  w fixes every point of F;
\item
  w fixes every point of \(\overline{\mathbf{F}}\) ;
\item
  w belongs to the subgroup of W generated by the reflections with respect to the walls of C containing F.
\end{enumerate}

For every subset A of E, denote by \(W(A)\) the subgroup of W consisting of the elements that fix every point of A.

PROPOSITION 2. Let A be a non-empty subset of E, let \(\mathfrak{H}_{\mathrm{A}}\) be the set of hyperplanes \(\mathrm{H} \in \mathfrak{H}\) containing A, let \(\mathrm{A}'\) be the intersection of the \(\mathrm{H} \in \mathfrak{H}_{\mathrm{A}}\) , and let F be a facet of E open in \(\mathrm{A}'\) (§1, no. 3, Prop. 7). Then \(\mathrm{W(A)} = \mathrm{W(A')} = \mathrm{W(F)}\) , and \(\mathrm{W(A)}\) is generated by the reflections with respect to the hyperplanes in \(\mathfrak{H}_{\mathrm{A}}\) .

We first prove the following assertion:

\begin{enumerate}
\def\labelenumi{(\Roman{enumi})}
\tightlist
\item
  Let C be a chamber, let x and y be two points of \(\overline{C}\) and let \(w \in W\) be such that \(w(x) = y\) . Then x = y and w belongs to the subgroup \(W_{N}\) , where N is the set of walls of C containing x.
\end{enumerate}

We argue by induction on the length \(q\) of \(w\) (relative to the set \(S\) of reflections with respect to the walls of \(C\) ), the case \(q = 0\) being obvious. If \(q \geqslant 1\) , there exist a wall \(H\) of \(C\) and an element \(w' \in W\) such that \(w = s_H w'\) and \(l(w') = q - 1\) . Since \(l(s_H w) < l(w)\) , the chambers \(C\) and \(w(C)\) are on opposite sides of \(H\) by Th. 1 of no. 2. Thus, \(\overline{C} \cap w(\overline{C}) \subset H\) , so \(y \in H\) . Thus \(y = w'(x)\) and the induction hypothesis implies that \(x = y\) and \(w' \in W_{\mathfrak{N}}\) . Since \(y \in H\) , it follows that \(H \in \mathfrak{N}\) , and hence that \(w = s_H w' \in W_{\mathfrak{N}}\) , completing the proof of (I).

Proof of Theorem 2: this follows from (I) and Lemma 2.

Proof of Proposition 1: we know that two distinct facets are disjoint and have distinct closures (§1, no. 2, Cor. of Prop. 3). The equivalence of (i), (ii) and (iii) follows. On the other hand, it is clear that

\[
(\mathrm{vii}) \Longrightarrow (\mathrm{vi}) \Longrightarrow (\mathrm{v}) \Longrightarrow (\mathrm{iv}) \Longrightarrow (\mathrm{i})
\]

and assertion (I) shows that (i) \(\Longrightarrow\) (vii).

Proof of Proposition 2: let \(\mathbf{A}''\) be the affine subspace of \(\mathbf{E}\) generated by A. Clearly, \(\mathrm{W(A)} = \mathrm{W(A'')}\) . By Prop. 7 of §1, no. 3, there exists a point \(x \in \mathbf{A}''\) that does not belong to any hyperplane \(\mathbf{H} \in \mathfrak{H} - \mathfrak{H}_{\mathbf{A}}\) . Let \(\mathbf{F}_x\) be the facet containing \(x\) : it is open in \(\mathbf{A}''\) and Prop. 1 shows that

\[
\mathrm{W} \left(\mathrm{F} _ {x}\right) \subset \mathrm{W} \left(\mathrm{A} ^ {\prime}\right) \subset \mathrm{W} (\mathrm{A}) = \mathrm{W} \left(\mathrm{A} ^ {\prime \prime}\right) \subset \mathrm{W} (\{x \}) = \mathrm{W} \left(\mathrm{F} _ {x}\right),
\]

hence \(\mathrm{W(A) = W(A') = W(F_x)}\) . Replacing A by F, we also have

\[
\mathrm{W} (\mathrm{A}) = \mathrm{W} (\mathrm{F}),
\]

hence the proposition.

Remarks. 1) It follows from Th. 2 that, if C is a chamber and F a facet, there exists a unique facet \(F'\) contained in \(\overline{C}\) that is transformed into F by an element of W.

\begin{enumerate}
\def\labelenumi{\arabic{enumi})}
\setcounter{enumi}{1}
\item
  It follows from Props. 1 and 2 that, for any non-empty subset A of E, there exists a point \(a \in E\) such that \(W(A) = W(\{a\})\) ; moreover, the group \(W(A)\) is a Coxeter group (Th. 1).
\item
  Let \(\mathbf{C}\) be a chamber of \(\mathbf{E}\) and \(\mathbf{S}\) the set of reflections with respect to the walls of \(\mathbf{C}\) . Let \(w \in W\) and let \((s_1, \ldots, s_q)\) be a reduced decomposition of \(w\) with respect to \(\mathbf{S}\) . If \(x \in \overline{\mathbf{C}}\) is fixed by \(w\) , then \(s_j(x) = x\) for all \(j\) : this follows from Prop. 1 above and Cor. 1 of Prop. 7 of Chap. IV, §1.
\end{enumerate}

